\section{Discussion}
\subsection{Interpretation of results}

Using multi-agent reinforcement learning simulations to investigate the role of two empathy components (cognitive and motivational) in promoting pro-environmental behaviors, we hypothesized that the sustainability of the system dynamics, measured by episode length, would be enhanced by both components of empathy. Furthermore, along with Du et al. \cite{du_2024}, we expected fairness (equality in resource acquisition) to be favored by empathy.

\subsubsection{Empathy and Sustainability}
In line with (TH1), Study 1 revealed significant main effects of the presence of both cognitive and motivational empathy on episode length, as well as a significant, sub-additive interaction between the two components. Motivational empathy emerged as the primary driver of sustainability: simulations in which agents received an empathic reward were associated with substantially longer episodes before depletion, regardless of access to affective signals. The small magnitude of the difference between the Blind/Empathic and Full Empathy conditions suggests that motivational empathy alone may be sufficient to sustain the resource at a comparable level, and that cognitive access to others' emotional states only marginally amplifies this effect under simple binary action selection. Additionally, in line with (TH3), Study 2 confirmed this dose-response relationship, with a strong positive correlation between the degree of altruistic motivation and resource sustainability.

These results are consistent with the hypothesis of empathy-altruism \cite{batson1997perspective,Batson01041991,batson2022empathy}, which proposes that empathic concern, an emotional response oriented to other people that is congruent with the perceived welfare of another, is the primary driver of prosocial behavior. As a process, empathic concern is triggered by two conditions: a belief that another is in need, and an intrinsic valuation of that other's welfare. Once elicited, it generates a motivational state in which egoistic and altruistic motivations interact through a hedonic calculus, ultimately shaping action selection. Our framework directly instantiates this process: the valence signal broadcast by each agent is congruent with its consumption satisfaction, and the empathic reward received by peers is a direct function of this signal. When others signal distress (negative valence), agents with a nonzero empathic reward coefficient are penalized for consuming, functionally replicating the trigger of empathic concern. The empathic weight itself models the hedonic/utility calculus at the heart of the hypothesis, interpolating between purely egoistic (altruistic weight $\alpha_e = 0$) and predominantly altruistic ($\alpha_e \rightarrow 1$) motivational states.

This operationalization also maps onto the distinction between reactive empathy and perspective-taking. Reactive empathy, feeling for another by experiencing an emotion of congruent valence \cite{batson1997perspective}, is captured by the motivational component: agents experience a reward signal of the same valence as their peers' affective states, increasing their empathic concern. Perspective-taking, by contrast, improves the adjustment to a more complex state space. In our framework, access to others' affective signals (cognitive empathy condition) provides agents with the perceptual input needed for such estimation. This may explain why the interaction between cognitive and motivational empathy on sustainability is sub-additive: perspective-taking refines the expression of altruistic motivation but is not itself a source of that motivation. Together, this aligns with dual-interest theory \cite{lynneTheoreticalFoundationEmpathy2016} and cognitive-affective-motivational conceptualizations of empathy \cite{decetyNeuralPathwaysDevelopment2015,dewaalMammalianEmpathyBehavioural2017}. That said, depletion remained nearly systematic across all conditions, suggesting that empathy as operationalized here delays but does not prevent the Tragedy of the Commons. This points to the need for complementary mechanisms, such as internal representations of resource state or institutional constraints, to achieve genuine sustainability \cite{ostromGeneralFrameworkAnalyzing2009,vanlangeHowWeCan2008}.

\subsubsection{Empathy and Fairness}
Contrary to (TH2), both cognitive and motivational empathy were associated with higher Gini coefficients, indicating greater inequality in resource distribution. Study 2 further confirmed this pattern, with an opposite result to (TH4) with altruistic motivation strongly and positively correlated with inequality.

At first glance, this contradicts the empathy-altruism hypothesis, which would predict that other-regarding motivation promotes equitable as well as sustainable outcomes. However, the hypothesis specifies that empathic concern targets welfare, not equality \cite{Batson01041991}. The discrete set of emotions constituting empathic concern is oriented toward the perceived need of another, not toward distributive justice between all others. In a shared resource environment, an agent experiencing empathic concern may consistently restrain its own consumption to benefit peers, but this self-restraint does not compel equal distribution across all agents. Rather, as agents follow heterogeneous learning trajectories under identical algorithms, differential consumption rates can emerge, producing objective inequality as an unintended byproduct of other-regarding motivation.

The reward structure reinforces this dynamic. Agents who consume less still achieve comparable aggregate rewards because the social component of the reward compensates for reduced personal gain. This reduces sensitivity to self-disadvantageous inequity: since perceived aggregate utility remains stable regardless of consumption asymmetries, there is less pressure to equalize access \cite{heInfluenceEmpathicConcern2022}. Empathic concern may thus decouple sustainability from fairness.

The significant interaction observed in Study 1 offers a partial qualification: the inequality associated with motivational empathy was attenuated when agents could perceive others' affective states. Within the empathy-altruism framework, this is interpretable through the role of perspective-taking: access to others' valence signals improves the accuracy of need estimation \cite{batson1997perspective}, allowing agents to direct their care more precisely and reducing the polarization of consumption patterns. Without affective signals, reactive empathy remains diffuse and undirected; with them, agents can modulate restraint in response to specific others' perceived states, producing a more needs-responsive and comparatively less unequal distribution \cite{moreno-casasEthicsCareTragedy2021}.

These results raise broader theoretical questions: to what extent can empathy offset personal disadvantage? Are self-interest and social rewards integrated into a unified utility signal or processed as separate domains? And how would these dynamics change if the resource were essential for survival, altering the valence and intensity of the empathic concern elicited? Developing precise computational models with categorical distinctions between empathy, normative influence, and instrumental reasoning remains essential to advance our understanding of pro-environmental behavior.


\subsection{Limitations and future work}
\subsubsection{Limitations}
Simulation parameters were selected through brute-force exploration, following an approach inspired by Perolat et al.\ (2017) \cite{perolatMultiagentReinforcementLearning2017}. The simulations were carried out under conditions where a sufficient number of agents was known to reliably produce resource depletion, allowing the observation of emergent behaviors under scarcity.

Several features of the current design limit direct transferability to human behavior. First, the affective signaling mechanism relies on a linear mapping from consumption rate to valence, which may not capture the nonlinear dynamics of human emotional experience. Alternative mappings,  such as sigmoid transformations or bio-inspired computational models of emotion \cite{gannouniAdaptiveEmotionDetection2020,joffilyEmotionalValenceFreeenergy2013},  could yield richer affective signals and increase the plausibility of our framework as a model of human empathic processing. Second, agents in the current simulations discount future rewards at a rate of $\gamma = 0.99$, meaning they weigh near-future outcomes nearly as heavily as immediate ones. This contrasts with human temporal discounting, where future rewards are systematically devalued, an effect that intensifies under uncertainty \cite{storyDiscountingFutureReward2024}. Testing lower discount factors would improve ecological validity by assessing whether the sustainability benefits of empathy hold under more human-like temporal preferences. Third, research suggests that empathic valuation and behavioral investment in others' needs decay over time \cite{colemanIntertemporalEmpathyDecline2024}, a dynamic not currently captured in our framework. Fourth, even resource consumption in real-world contexts is rarely binary; human decisions involve graded, context-sensitive trade-offs that a two-action model cannot fully represent.

\subsubsection{Future Directions}
Beyond addressing these limitations, several extensions are planned. Moving to a 2D spatial environment would allow the study of proximity-based empathic dynamics and more realistic foraging behaviors. Exploring alternative reward functions, including non-linear formulations and need-sensitive consumption penalties, would allow a more fine-grained test of the conditions under which empathy promotes both sustainability and fairness. Finally, fitting these simulations to data from analogous human behavioral tasks could provide parameter estimates with greater empirical grounding, bridging the gap between computational modeling and psychological theory.
While this work provides insights into the role of empathy in pro-environmental behaviors, it is not sufficient. The transferability of our findings to human and other contexts is limited.

\subsection{Conclusion}
This study presents a novel multi-agent reinforcement learning framework for investigating the computational role of empathy in common-pool resource dilemmas. By operationalizing empathy through affective signaling observation (cognitive component) and reward shaping (motivational component), we observed that both components contribute independently and interactively to resource sustainability. Motivational empathy emerged as the primary driver of reduced overconsumption, while cognitive empathy modulated the fairness of resource distribution. Contrary to our initial hypotheses, higher empathic motivation was associated with greater inequality, an effect partially mitigated by access to peers' affective states.

These findings provide initial computational grounding for empathy-based theories of pro-environmental behavior and confirm that empathy can be a relevant mechanism for cooperative resource management beyond purely institutional or normative approaches. Future work should extend this framework to richer environments, including 2D spatial settings, continuous action spaces, and resources with regenerative dynamics, and to explore how these simulated results can be connected to human behavioral data through model fitting.
